% =============================================================================
% Q5 -- bench and simulated-use testing
% =============================================================================
\clearpage
\hypertarget{q5-overview}{}
\sectiontitle{Q5. Bench and Simulated-Use Testing}

\providecommand{\reportlink}[2]{%
  \hyperlink{#1}{\textcolor{AUBRed}{\bfseries\underline{#2}}}}

% Compact, non-floating literature panels used throughout Q5. A fixed image
% height keeps each multi-paper plate on one page while preserving aspect ratio.
\providecommand{\qfivelitpanel}[3]{%
  \begin{minipage}[t]{0.485\linewidth}
    \centering
    \includegraphics[width=\linewidth,height=#1,keepaspectratio]{#2}\par
    \vspace{0.18em}
    \raggedright{\fontsize{7.0}{8.05}\selectfont #3\par}
  \end{minipage}}

\qtwosubsection{5.1 Literature-backed test programme and evidence boundary}

{\fontsize{9.3}{10.95}\selectfont
Q5 converts the literature synthesis, functional blocks, and specifications
into a controlled test programme. The cited papers establish sensing
principles, reference apparatus, loading sequences, environmental ranges, and
reporting metrics. They do not establish the performance of the present
ink--printer--PDMS--sensor combination. The project must generate its own raw
records, calibration models, uncertainty estimates, plots, and pass/partial/
fail decisions. This structure follows the assignment requirement to state the
test, materials and tools, rationale, expected outcome, and impact on the
overall design.\par}

\evidencebox{Interpretation rule:}{An ``expected outcome'' below is a
predeclared engineering expectation, not a measured result. A literature value
is marked as precedent; a Q4 value is a project specification; only a recorded
project measurement can close a verification item.}

\vspace{0.45em}
\begingroup
\fontsize{7.55}{8.85}\selectfont
\setlength{\tabcolsep}{2.5pt}\renewcommand{\arraystretch}{1.05}
\begin{tabularx}{\linewidth}{>{\centering\arraybackslash\bfseries}p{1.25cm}Y >{\raggedright\arraybackslash}p{4.2cm} >{\raggedright\arraybackslash\bfseries\color{AUBRed}}p{2.5cm}}
\rowcolor{AUBRed}
\color{white}\textbf{Origin} & \color{white}\textbf{Clickable evidence chain} &
\color{white}\textbf{Q4 requirement family} & \color{white}\textbf{Q5 test IDs} \\
\reportlink{q1-question-1}{Q1-1} &
\reportlink{q2-functional-blocks}{Q2 2.1} $\rightarrow$
\reportlink{q3-functional-blocks}{Q3A} & GEO, MAT, PRT, SYS & F0, X1, SYS1 \\
\rowcolor{AUBPale}
\reportlink{q1-question-2}{Q1-2} &
\reportlink{q2-humidity}{Q2 2.4.1} $\rightarrow$
\reportlink{q3-microclimate-evidence}{Q3 3.3} & HUM, ENV, SAFE & H1, H2, D1 \\
\reportlink{q1-question-3}{Q1-3} &
\reportlink{q2-pressure-shear}{Q2 2.3} / \reportlink{q2-temperature}{2.4}
$\rightarrow$ \reportlink{q3-operational-definitions}{Q3 3.1} &
MECH, TEMP, HUM & M1--M3, T1, H1 \\
\rowcolor{AUBPale}
\reportlink{q1-question-4}{Q1-4} &
\reportlink{q2-current-build}{Q2 2.2} $\rightarrow$
\reportlink{q2-pdms-printing}{Q2 2.5} & MFG, PRT, MAT & F0, R1 \\
\reportlink{q1-question-5}{Q1-5} &
\reportlink{q2-whole-liner}{Q2 2.7} $\rightarrow$
\reportlink{q3-measurement-models}{Q3 3.5} & DAQ, MAP, ALG & X1, MAP1, SYS1 \\
\rowcolor{AUBPale}
\reportlink{q1-question-6}{Q1-6} &
\reportlink{q2-current-build}{Q2 2.2} $\rightarrow$
\reportlink{q3-metrics}{Q3 3.7} & MAP, GEO & MAP1, S1 \\
\reportlink{q1-question-7}{Q1-7} &
\reportlink{q2-pdms-printing}{Q2 2.5} $\rightarrow$
\reportlink{q3-validation-outputs}{Q3 3.8} & DUR, ENV, SAFE & D1, D2, R1 \\
\end{tabularx}
\endgroup
{\fontsize{7.45}{8.6}\selectfont\color{AUBGray}
\textbf{Table Q5-1. Clickable Q1--Q5 traceability.} The links open the detailed
literature, operational definitions, and metrics instead of repeating those
sections. The complete specification set is in
\reportlink{q4-requirement-basis}{Q4 Section 4.1}.\par}

\vspace{0.55em}
\begin{center}
\begin{tikzpicture}[>=stealth,node distance=2.5mm]
  \tikzstyle{qfivegate}=[draw=AUBLine,fill=AUBPale,rounded corners=1pt,
    text width=29mm,minimum height=10mm,align=center,
    font=\fontsize{6.7}{7.6}\selectfont]
  \node[qfivegate] (a) {F0\quad print, geometry, continuity};
  \node[qfivegate,right=of a] (b) {M/T/H\quad individual calibration};
  \node[qfivegate,right=of b] (c) {X/MAP\quad integration and location};
  \node[qfivegate,right=of c] (d) {D\quad environmental and cyclic};
  \node[qfivegate,right=of d] (e) {S/SYS\quad phantom and release};
  \draw[->,thick,color=AUBRed] (a)--(b);
  \draw[->,thick,color=AUBRed] (b)--(c);
  \draw[->,thick,color=AUBRed] (c)--(d);
  \draw[->,thick,color=AUBRed] (d)--(e);
\end{tikzpicture}
\end{center}
{\fontsize{7.5}{8.7}\selectfont\color{AUBGray}
\textbf{Figure Q5-1. Gated execution sequence.} A failure at an earlier gate
triggers rework or a documented partial result; it cannot be hidden by a later
combined score. This implements the sequence defined in
\reportlink{q3-execution-sequence}{Q3 Section 3.6} and
\reportlink{q4-design-gate}{Q4 Section 4.7}.\par}

% -----------------------------------------------------------------------------
\clearpage
\hypertarget{q5-release}{}
\sectiontitle{Q5. Specimen Release and Print-Quality Bench Tests}

\qtwosubsection{5.2 F0 --- release each printed specimen before calibration}

{\fontsize{9.15}{10.75}\selectfont
The PET sheet in \reportlink{q2-current-build}{Q2 Section 2.2} demonstrates
first-stage artwork transfer; it is not the calibration specimen. F0 applies
to direct-on-PDMS coupons and then to the complete 56\,$\times$\,55\,$\times$
5.2~mm slab. Liew \textit{et al.} used optical microscopy and design-to-print
dimensional comparison before electrical temperature characterization
\cite{liew2020}. Albrecht \textit{et al.} combined optical inspection,
profilometry, and electrical monitoring for printed electrodes and PDMS-based
conductors \cite{albrecht2019,albrecht2018stretch}. These methods support the
inspection categories below; the acceptance limits remain those declared in
\reportlink{q4-fabrication}{Q4 Section 4.3}.\par}

\begingroup
\fontsize{7.8}{9.15}\selectfont
\setlength{\tabcolsep}{2.7pt}\renewcommand{\arraystretch}{1.07}
\begin{tabularx}{\linewidth}{>{\raggedright\arraybackslash\bfseries\color{AUBRed}}p{1.05cm}Y Y Y}
\rowcolor{AUBRed}
\color{white}\textbf{Step} & \color{white}\textbf{Materials and tools} &
\color{white}\textbf{Recorded evidence and rationale} &
\color{white}\textbf{Expected outcome and design impact} \\
F0-1 & Printed coupon/slab, CAD export, caliper, optical microscope or USB
microscope, scale target. & Photograph every modality and pad; measure line
width, gap, active footprint, edge clearance, slab dimensions, and layer
thickness. This detects print spreading, nozzle loss, mold error, and
misregistration before electrical fitting. & Complete image set and dimension
table. Out-of-tolerance geometry changes the print file, surface treatment,
drop spacing, mold, or alignment method before calibration. \\
\rowcolor{AUBPale}
F0-2 & Four-wire meter where available; otherwise calibrated DMM, guarded LCR
meter, probes, documented cables. & Record conductor resistance, open/short
status, pad-to-pad isolation, and the zero-load baseline of every channel.
Repeat after 10~min to identify unstable contacts. & Every channel has a unique
ID and stable baseline. An open, short, or drifting contact is repaired or the
specimen is rejected; missing channels are not interpolated away. \\
F0-3 & Profilometer where available; alternatively a sacrificial step coupon,
microscope focus method, or cross-sectional witness. & Record printed height
and PDMS swelling/relief. The same-printer study reported a much larger
apparent profile on PDMS than PET, partly from solvent-induced swelling
\cite{albrecht2018stretch}. & Geometry is reported as measured, not inferred
from deposited ink volume. Excess swelling changes drying, surface treatment,
or the chosen ink route. \\
\rowcolor{AUBPale}
F0-4 & Adhesive-tape screen, gentle bend mandrel, insulation probe, camera. &
Inspect adhesion, cracking, exposed silver, encapsulation discontinuity, and
continuity before and after a non-destructive handling screen. & Only intact,
insulated, electrically continuous specimens advance. Failure changes
encapsulation, trace geometry, or handling fixtures. \\
F0-5 & Batch traveller, printer log, ink lot, PDMS ratio/cure record, specimen
label. & Preserve printer settings, substrate preparation time, cure, drying,
inspection, and operator. Direct-on-PDMS performance is sensitive to surface
treatment and process history \cite{q4abukhalaf2018,q4abukhalaf2019}. & A
traceable build record links every later result to one specimen and prevents
selective reporting of only successful prints. \\
\end{tabularx}
\endgroup

\vspace{0.5em}
\begin{center}
\begin{tikzpicture}[>=stealth,node distance=3mm]
  \tikzstyle{trav}=[draw=AUBLine,fill=AUBPale,rounded corners=1pt,
    text width=34mm,minimum height=8mm,align=center,
    font=\fontsize{6.8}{7.7}\selectfont]
  \node[trav] (a) {build record\\ink + PDMS + printer};
  \node[trav,right=of a] (b) {geometry\\image + dimensions};
  \node[trav,right=of b] (c) {electrical\\continuity + isolation};
  \node[trav,right=of c] (d) {release\\advance / rework / reject};
  \draw[->,color=AUBRed,thick] (a)--(b);
  \draw[->,color=AUBRed,thick] (b)--(c);
  \draw[->,color=AUBRed,thick] (c)--(d);
\end{tikzpicture}
\end{center}
{\fontsize{7.5}{8.7}\selectfont\color{AUBGray}
\textbf{Figure Q5-2. F0 specimen traveller.} The released specimen ID becomes
the key used by every calibration file, CAD coordinate, photograph, and
durability record.\par}

% -----------------------------------------------------------------------------
\clearpage
\hypertarget{q5-lit-fabrication}{}
\sectiontitle{Q5 Literature Evidence Plate: Direct Printing and Release}

{\fontsize{8.75}{10.25}\selectfont
These four source figures are the visual basis for F0. They show the substrate
effect, an implemented PDMS-printing workflow, a geometry/process response,
and the electrical calibration obtained from an inkjet-printed meander. They
support the methods in \reportlink{q5-release}{Q5 Section 5.2} and the material
route in \reportlink{q2-pdms-printing}{Q2 Section 2.5}; they are not results
from the present 5.2~mm slab.\par}

\vspace{0.55em}
\noindent
\qfivelitpanel{0.195\textheight}
  {paper_sources/figures_selected/albrecht2018_fig2_pdms_pet_height_profile.png}
  {\textbf{(a) Substrate-dependent printed profile.} Original Fig.~2 in
  Albrecht \textit{et al.}: silver printed with the same settings produced an
  apparent mean height of 3.8~$\mu$m on PDMS and 0.6~$\mu$m on PET, with PDMS
  swelling identified as a major contributor \cite{albrecht2018stretch}.
  \href{https://doi.org/10.3390/polym10121413}{Open DOI}.}
\hfill
\qfivelitpanel{0.195\textheight}
  {paper_sources/q5_figures/crops/abukhalaf2018_fig4.png}
  {\textbf{(b) Implemented direct-on-PDMS process.} Original Fig.~4 in
  Abu-Khalaf \textit{et al.}: pre-stretching and clamping, surface treatment,
  inkjet printing, and oven sintering of silver traces on PDMS
  \cite{q4abukhalaf2018}.
  \href{https://doi.org/10.3390/ma11122377}{Open DOI}.}

\vspace{0.75em}
\noindent
\qfivelitpanel{0.185\textheight}
  {paper_sources/q5_figures/crops/abukhalaf2019_fig4.png}
  {\textbf{(c) Geometry--process trade-off.} Original Fig.~4 in Abu-Khalaf
  \textit{et al.}: response-surface prediction of breakdown strain versus
  printed line width and number of layers. It supports recording geometry and
  print history before durability interpretation \cite{q4abukhalaf2019}.
  \href{https://doi.org/10.3390/ma12203329}{Open DOI}.}
\hfill
\qfivelitpanel{0.185\textheight}
  {paper_sources/q5_figures/crops/liew2020_fig5.png}
  {\textbf{(d) Printed-meander calibration precedent.} Original Fig.~5 in
  Liew \textit{et al.}: six resistance--temperature series for commercial and
  in-house AgNP inks. F0 first verifies the printed geometry and continuity;
  T1 then fits this type of calibration \cite{liew2020}.
  \href{https://doi.org/10.3390/ecsa-7-08216}{Open DOI}.}

\vfill
{\fontsize{6.7}{7.7}\selectfont\color{AUBGray}
\textbf{Figure Q5-L1. Literature evidence for specimen release and direct
printing.} Panels are reproduced as source excerpts with their original
figure numbers and citations. Open-access MDPI figures are used under their
article licences. The acceptance thresholds remain the project requirements
in \reportlink{q4-fabrication}{Q4 Section 4.3}.\par}

% -----------------------------------------------------------------------------
\clearpage
\hypertarget{q5-mechanical}{}
\sectiontitle{Q5. Pressure and Directional-Shear Bench Testing}

\qtwosubsection{5.3 M1--M3 --- calibration, separation, and spatial coverage}

{\fontsize{8.9}{10.45}\selectfont
The closest direct force-cell precedent used a voice-coil actuator with an
integrated force sensor, an Agilent E4980A LCR meter at 100~kHz and 1~V AC, and
0.1--8~N loading in 20\% increments every 5~s, followed by unloading and three
repeated cycles \cite{albrecht2018conf,albrecht2019}. Q5 reuses the stepped
loading, unloading, directional reversal, and independent force reference. It
does not copy the paper's calibration coefficients because the present sensor
is printed and enclosed in a different 5.2~mm PDMS structure. The project
envelope and gates are clickable in
\reportlink{q4-mechanical}{Q4 Section 4.4}.\par}

\begin{center}
\begin{tikzpicture}[x=1cm,y=1cm,>=stealth]
  \fill[AUBPale] (0,0) rectangle (13.6,1.0);
  \draw[fill=white,draw=AUBLine] (4.8,1.0) rectangle (8.8,1.45);
  \node[font=\fontsize{6.8}{7.6}\selectfont] at (6.8,1.22) {released 5.2 mm slab on rigid fixture};
  \draw[fill=AUBGray!20,draw=AUBGray] (6.25,1.45) rectangle (7.35,1.75);
  \node[font=\fontsize{6.2}{7}\selectfont] at (6.8,1.60) {known footprint};
  \draw[->,very thick,color=AUBRed] (6.8,2.8)--(6.8,1.78);
  \node[font=\fontsize{7}{8}\selectfont\bfseries,text=AUBRed] at (7.55,2.5) {$F_N$};
  \draw[<->,very thick,color=AUBBlue] (4.8,1.95)--(8.8,1.95);
  \node[font=\fontsize{7}{8}\selectfont\bfseries,text=AUBBlue] at (9.25,1.95) {$\pm F_T$};
  \node[font=\fontsize{6.5}{7.3}\selectfont,align=center] at (2.2,1.25)
    {load cell or calibrated force gauge\\plus vertical/tangential stage};
  \node[font=\fontsize{6.5}{7.3}\selectfont,align=center] at (11.4,1.25)
    {LCR/DAQ records $C_1$--$C_4$\\and a common time base};
  \draw[->,color=AUBGray] (3.45,1.25)--(4.65,1.25);
  \draw[->,color=AUBGray] (8.95,1.25)--(10.1,1.25);
\end{tikzpicture}
\end{center}
{\fontsize{7.5}{8.7}\selectfont\color{AUBGray}
\textbf{Figure Q5-3. Mechanical bench arrangement.} A motorized actuator is
preferred for controlled ramps; calibrated masses, a force gauge, and a
low-friction translation stage are acceptable lower-cost implementations if
force and displacement are recorded independently.\par}

\vspace{0.35em}
\begingroup
\fontsize{7.65}{8.95}\selectfont
\setlength{\tabcolsep}{2.6pt}\renewcommand{\arraystretch}{1.04}
\begin{tabularx}{\linewidth}{>{\raggedright\arraybackslash\bfseries\color{AUBRed}}p{1.0cm}Y Y Y}
\rowcolor{AUBRed}
\color{white}\textbf{Test} & \color{white}\textbf{Procedure and tools} &
\color{white}\textbf{Rationale and expected outcome} &
\color{white}\textbf{Output and design impact} \\
M1 & Apply increasing/decreasing normal pressure from 0 to 250~kPa through a
measured footprint; use at least three cycles per level and held-out points.
Record force, displacement, $C_1$--$C_4$, dwell, and recovery. & Establishes
the usable range, saturation, viscoelastic hysteresis, and baseline recovery.
All four capacitors should change coherently under normal loading. & Calibration
curve, residuals, $R^2$, NRMSE, hysteresis, coefficient of variation, and
recovery. Failure changes footprint, spacer/cover thickness, preload,
electrode overlap, or fitting model. \\
\rowcolor{AUBPale}
M2 & At three or more normal preloads, apply signed shear from $-50$ to
$+50$~kPa in both slab axes; reverse direction and include 5, 30, and 300~s
dwells. Record an independent tangential force/displacement reference. & Tests
direction separation and normal-load cross-coupling. Opposed capacitor pairs
should change with opposite signs for the driven axis, while the non-target
pair remains bounded \cite{albrecht2019}. & Vector sensitivity, angular error,
cross-axis response, hysteresis, and recovery. Failure changes alignment,
four-sector geometry, dielectric thickness, or compensation matrix. \\
M3 & Apply M1/M2 stimuli at registered sensor centres, between adjacent cells,
and near slab boundaries. Randomize positions and blind the software operator.
& Determines whether spacing and placement detect localized conditions rather
than only centre loads. & Coverage mask, missed-event rate, true/detected
coordinate pairs, and edge error. Failure changes sensor density, placement,
interpolation boundaries, or the claimed coverage region. \\
\end{tabularx}
\endgroup

\vspace{0.35em}
\begin{align*}
P&=F_N/A, & \tau&=F_T/A, &
C_z&=\tfrac14(C_1+C_2+C_3+C_4),\\[-0.3em]
C_x&=C_1-C_3, & C_y&=C_2-C_4, &
H_{FS}&=\max|y_{\uparrow}-y_{\downarrow}|/FS.
\end{align*}
{\fontsize{7.45}{8.6}\selectfont\color{AUBGray}
Pressure and shear are assigned only after the reference footprint, force,
sign convention, calibration version, and uncertainty are recorded. A force
signal alone is not proof of slip; relative motion requires the independent
reference described in \reportlink{q3-ground-truth}{Q3 Section 3.4}
\cite{q3tang2022,noll2018}.\par}

% -----------------------------------------------------------------------------
\clearpage
\hypertarget{q5-lit-mechanical}{}
\sectiontitle{Q5 Literature Evidence Plate: Force and Motion Benches}

{\fontsize{8.75}{10.25}\selectfont
The closest published apparatus separate controlled normal force, controlled
shear, and relative limb--socket motion. This plate therefore supports the
independent references and signed loading sequence in
\reportlink{q5-mechanical}{Q5 Section 5.3}, not a claim that the present slab
will reproduce the published transfer functions.\par}

\vspace{0.55em}
\noindent
\qfivelitpanel{0.19\textheight}
  {paper_sources/q5_figures/crops/albrecht2018conf_fig1.png}
  {\textbf{(a) Controlled normal/shear fixtures.} Original Fig.~1 in Albrecht
  \textit{et al.} (ALLSENSORS 2018): voice-coil/force-sensor arrangements for
  normal loading and rotatable signed shear \cite{albrecht2018conf}. The image
  directly motivates separate $F_N$ and $F_T$ references in M1--M2.}
\hfill
\qfivelitpanel{0.19\textheight}
  {paper_sources/q5_figures/crops/albrecht2019_fig9.png}
  {\textbf{(b) Normal/shear channel separation.} Original Fig.~9 in
  Albrecht \textit{et al.}: time-domain normal-force response and opposite
  capacitance signs under reversed shear direction in the printed sensor
  \cite{albrecht2019}.
  \href{https://doi.org/10.1155/2019/1864239}{Open DOI}.}

\vspace{0.75em}
\noindent
\qfivelitpanel{0.205\textheight}
  {paper_sources/q5_figures/crops/tang2022_fig1.png}
  {\textbf{(c) Synchronized interface stress and motion.} Original Fig.~1 in
  Tang \textit{et al.}: pressure/shear acquisition synchronized with a 3D
  motion-capture system and force plate \cite{q3tang2022}. This is the
  precedent for treating slip or pistoning as an independently labelled event.
  \href{https://doi.org/10.1177/09544119221110712}{Open DOI}.}
\hfill
\qfivelitpanel{0.205\textheight}
  {paper_sources/q5_figures/crops/noll2018_fig2.png}
  {\textbf{(d) Low-motion test rig and liner surfaces.} Original Fig.~2 in
  Noll \textit{et al.}: a movable cantilever, fixed 2D-motion sensor, and five
  liner/reference textures \cite{noll2018}. It supports an independent motion
  check under controlled surface and direction changes.
  \href{https://doi.org/10.3390/s18072170}{Open DOI}.}

\vfill
{\fontsize{6.7}{7.7}\selectfont\color{AUBGray}
\textbf{Figure Q5-L2. Mechanical validation precedents.} The two Albrecht
papers establish the force-cell apparatus and printed-capacitor response;
Tang and Noll establish the separate relative-motion reference required before
using the word ``slip.'' Rights remain with the cited authors/publishers where
an article licence is not explicitly stated in the supplied source.\par}

% -----------------------------------------------------------------------------
\clearpage
\hypertarget{q5-microclimate}{}
\sectiontitle{Q5. Temperature, Humidity, and Sweat-Protection Tests}

\qtwosubsection{5.4 T1, H1, and H2 --- separate thermal, vapour, and liquid evidence}

{\fontsize{8.75}{10.3}\selectfont
Liew \textit{et al.} inspected the printed meander and calibrated resistance
from 30 to 100~$^{\circ}$C in 10~$^{\circ}$C increments with a digital
multimeter \cite{liew2020}. J\"ager \textit{et al.} later characterized
inkjet-printed resistive sensors through temperature cycles and reported TCR,
hysteresis, non-repeatability, drift, and maximum error following an
IEC-based method \cite{jager2022}. A\"{\i}t-Mammar \textit{et al.} swept
20--90\%RH at 25~$^{\circ}$C in a climatic chamber, used increasing and
decreasing steps, and separately applied 20--80\%RH transients while measuring
with a reference hygrometer \cite{aitmammar2020}. Q5 adopts these measurement
patterns within the prosthetic-interface ranges specified in
\reportlink{q4-temp-humidity}{Q4 Section 4.5}; it keeps vapour calibration and
liquid-sweat protection as different tests.\par}

\begin{center}
\begin{tikzpicture}[>=stealth,node distance=5mm]
  \tikzstyle{envbox}=[draw=AUBLine,fill=AUBPale,rounded corners=1pt,
    text width=54mm,minimum height=22mm,align=center,
    font=\fontsize{6.8}{7.7}\selectfont]
  \node[envbox] (t) {\textbf{T1: temperature chamber/block}\\
    25--45~$^{\circ}$C, 5~$^{\circ}$C steps\\reference RTD/thermocouple + DMM/DAQ};
  \node[envbox,right=of t] (h) {\textbf{H1: vapour chamber}\\
    20--95\%RH, up/down steps\\reference hygrometer + LCR/DAQ};
  \node[envbox,right=of h] (s) {\textbf{H2: protected liquid exposure}\\
    controlled artificial-sweat volume\\barrier inspection + recovery calibration};
  \draw[->,color=AUBRed,thick] (t)--(h);
  \draw[->,color=AUBRed,thick] (h)--(s);
\end{tikzpicture}
\end{center}
{\fontsize{7.45}{8.6}\selectfont\color{AUBGray}
\textbf{Figure Q5-4. Ordered microclimate evidence.} T1 and H1 establish dry
calibration before H2 challenges the selected protective structure. Saturated
salt solutions provide literature-defined RH fixed points when a programmable
chamber is unavailable \cite{greenspan1977}; a calibrated reference hygrometer
must still be logged beside the specimen.\par}

\vspace{0.35em}
\begingroup
\fontsize{7.55}{8.85}\selectfont
\setlength{\tabcolsep}{2.5pt}\renewcommand{\arraystretch}{1.04}
\begin{tabularx}{\linewidth}{>{\raggedright\arraybackslash\bfseries\color{AUBRed}}p{0.95cm}Y Y Y}
\rowcolor{AUBRed}
\color{white}\textbf{Test} & \color{white}\textbf{Materials, tools, and sequence} &
\color{white}\textbf{Expected evidence} & \color{white}\textbf{Design impact} \\
T1 & Released slab, temperature chamber/block, co-located calibrated RTD or
thermocouple, DMM/DAQ. Apply 25--45~$^{\circ}$C at 5~$^{\circ}$C steps up and
down; stabilize, hold out intermediate temperatures, and repeat after bending.
& Resistance-temperature model, TCR, sensitivity, bias, RMSE, $t_{90}$,
hysteresis, non-repeatability, and drift. Expected: monotonic repeatable
response with Q4 error and recovery gates met. & Revise meander geometry,
readout current, thermal isolation, encapsulation, compensation, or sampling
rate. A wide 30--100~$^{\circ}$C coupon screen is not substituted for the
25--45~$^{\circ}$C integrated-slab calibration. \\
\rowcolor{AUBPale}
H1 & Vapour chamber or sealed salt-reference vessels, calibrated hygrometer,
temperature reference, LCR/DAQ. Apply 20--95\%RH up/down at two or more
temperatures, with held-out RH levels and dry recovery. & Capacitance-RH model,
temperature compensation, hysteresis, repeatability, $t_{10\rightarrow90}$,
$t_{90\rightarrow10}$, saturation, and dry baseline. Expected: stable vapour
response through the selected vent/barrier. & Revise IDE spacing/coating,
barrier area, vent geometry, temperature compensation, equilibration time, or
the stated RH range. \\
H2 & Artificial sweat of documented composition, micropipette or syringe,
hydrophobic vapour-permeable barrier, absorbent containment, camera, leakage
inspection, and H1 references. Apply controlled volumes outside the protected
region; dry and repeat held-out H1 points. & Liquid ingress, leakage current,
short/open status, peak saturation, recovery time, and post-exposure
calibration shift. Expected: vapour access without liquid contact with silver
or joints. & Reject the barrier, increase sealed perimeter, move the sensor,
separate the liquid state in software, or add a replaceable protective layer.
Do not interpret saturation as higher discomfort. \\
\end{tabularx}
\endgroup

\vspace{0.3em}
\begin{align*}
\alpha&=\frac{R(T)-R_0}{R_0(T-T_0)}, &
S_T&=\frac{\Delta R}{\Delta T}, &
D_{post}&=\frac{y_{post}-y_{pre}}{FS}\times100\%.
\end{align*}
{\fontsize{7.4}{8.55}\selectfont\color{AUBGray}
The published humidity response below 100~s and recovery below 5~s belong to
the CAB-coated Kapton device, not to the protected PDMS slab
\cite{aitmammar2020}. The barrier and PDMS enclosure require new project
response-time measurements.\par}

% -----------------------------------------------------------------------------
\clearpage
\hypertarget{q5-lit-microclimate}{}
\sectiontitle{Q5 Literature Evidence Plate: Temperature and Humidity}

{\fontsize{8.75}{10.25}\selectfont
These figures show the actual curve shapes and reference values behind T1 and
H1. They justify regression, error, hysteresis, and fixed-RH checks in
\reportlink{q5-microclimate}{Q5 Section 5.4}; H2 remains a separate liquid-sweat
protection test because neither dry calibration plot proves liquid exclusion.\par}

\vspace{0.55em}
\noindent
\qfivelitpanel{0.205\textheight}
  {paper_sources/q5_figures/crops/jager2022_fig5.png}
  {\textbf{(a) Standardized temperature error reporting.} Original Fig.~5 in
  J\"ager \textit{et al.}: a fitted characteristic curve paired with the
  cycle-by-cycle measurement-error curve. The same data are used to quantify
  TCR, hysteresis, non-repeatability, and maximum error \cite{jager2022}.
  \href{https://doi.org/10.3390/s22218145}{Open DOI}.}
\hfill
\qfivelitpanel{0.205\textheight}
  {paper_sources/q5_figures/crops/aitmammar2020_fig4.png}
  {\textbf{(b) Humidity upsweep/downsweep.} Original Fig.~4 in
  A\"{\i}t-Mammar \textit{et al.}: normalized capacitance versus RH for three
  devices, including the separation between adsorption and desorption curves
  \cite{aitmammar2020}. This is the visual precedent for the H1 hysteresis
  metric. \href{https://doi.org/10.1557/adv.2020.86}{Open DOI}.}

\vspace{0.8em}
\begin{center}
\includegraphics[width=0.90\linewidth,height=0.285\textheight,keepaspectratio]
  {paper_sources/q5_figures/crops/greenspan1977_table2.png}
\end{center}
{\fontsize{7.0}{8.05}\selectfont
\textbf{(c) Chamber-independent RH reference values.} Table~2 in Greenspan
lists equilibrium RH versus temperature for saturated salt solutions. Q5 may
use selected salts as fixed points only when the specimen temperature and a
calibrated reference hygrometer are recorded \cite{greenspan1977}.
\href{https://nvlpubs.nist.gov/nistpubs/jres/81A/jresv81An1p89_A1b.pdf}
{Open NIST full text}.\par}

\vfill
{\fontsize{6.7}{7.7}\selectfont\color{AUBGray}
\textbf{Figure Q5-L3. Microclimate calibration precedents.} The J\"ager and
A\"{\i}t-Mammar curves define outputs to calculate; Greenspan supplies
reference states, not sensor performance. All three must be regenerated for
the final protected PDMS slab using project data.\par}

% -----------------------------------------------------------------------------
\clearpage
\hypertarget{q5-integration-mapping}{}
\sectiontitle{Q5. Multimodal Cross-Sensitivity and Map Validation}

\qtwosubsection{5.5 X1 and MAP1 --- prove that a correct map follows correct channels}

{\fontsize{8.95}{10.55}\selectfont
Cross-sensitivity is tested by applying one controlled input while recording
every channel, followed by selected combined conditions. Kwak \textit{et al.}
demonstrated synchronized multimodal prosthetic-interface acquisition at
30~Hz, Patern\`o \textit{et al.} registered sensor locations and logged slower
temperature/RH data, and Turner \textit{et al.} compared spatial interpolation
on a three-dimensional socket \cite{q3kwak2020,paterno2020,q3turner2021}.
These precedents support synchronized acquisition and spatial visualization;
they do not prove this slab's coordinate mapping or combined score. The data
schema and claim boundary are in \reportlink{q4-mapping}{Q4 Section 4.6}.\par}

\begingroup
\fontsize{7.65}{8.95}\selectfont
\setlength{\tabcolsep}{2.6pt}\renewcommand{\arraystretch}{1.06}
\begin{tabularx}{\linewidth}{>{\raggedright\arraybackslash\bfseries\color{AUBRed}}p{1.05cm}Y Y Y}
\rowcolor{AUBRed}
\color{white}\textbf{Run} & \color{white}\textbf{Controlled input} &
\color{white}\textbf{Channels and ground truth} &
\color{white}\textbf{Required interpretation} \\
X1-1 & Normal pressure only, several levels and coordinates. & Force/load cell,
$C_1$--$C_4$, temperature, RH, timestamps. & Quantify mechanical target
response and apparent thermal/RH interference; do not remove it by averaging. \\
\rowcolor{AUBPale}
X1-2 & Signed shear only under fixed normal preloads. & Tangential force and
displacement, four capacitances, temperature, RH. & Quantify target-axis,
cross-axis, and non-mechanical response. \\
X1-3 & Temperature only at controlled RH. & Reference temperature, meander
resistance, force channels, RH channel. & Separate temperature coefficient from
thermomechanical baseline shifts. \\
\rowcolor{AUBPale}
X1-4 & Vapour RH only at controlled temperature, then protected liquid sweat. &
Reference RH/temperature, IDE capacitance, leakage/continuity, all other
channels. & Separate vapour response, swelling/load redistribution, liquid
saturation, and electrical failure. \\
X1-5 & Predeclared combinations: pressure+shear; temperature+RH;
pressure+shear+heat+moisture. & All references and channels on one monotonic
clock; raw and filtered records retained. & Fit compensation only on training
runs; verify on held-out combinations. Report each physical-unit map before
any interface-condition score. \\
\end{tabularx}
\endgroup

\vspace{0.4em}
\begin{minipage}[t]{0.52\textwidth}
  \vspace{0pt}\centering
  \begin{tikzpicture}[x=0.72cm,y=0.72cm,>=stealth]
    \draw[step=1,AUBLine] (0,0) grid (8,5);
    \foreach \x/\y in {1/1,3/1,5/1,7/1,1/3,3/3,5/3,7/3}{
      \fill[AUBBlue] (\x,\y) circle (2pt);}
    \fill[AUBRed] (5.6,2.6) circle (3.2pt);
    \draw[AUBRed,thick] (5.6,2.6) circle (7pt);
    \fill[AUBYellow!80!black] (5.1,2.9) circle (2.7pt);
    \draw[->,thick,color=AUBRed] (5.6,2.6)--(5.1,2.9);
    \node[font=\fontsize{6.4}{7.2}\selectfont,text=AUBRed,anchor=west]
      at (5.8,2.35) {true input};
    \node[font=\fontsize{6.4}{7.2}\selectfont,text=AUBYellow!70!black,anchor=east]
      at (4.85,3.2) {detected};
    \node[font=\fontsize{6.2}{7}\selectfont,text=AUBGray] at (4,-0.55)
      {CAD coordinate plane with measured sensor sites};
  \end{tikzpicture}
\end{minipage}\hfill
\begin{minipage}[t]{0.45\textwidth}
  \vspace{0pt}
  {\fontsize{8.25}{9.7}\selectfont
  For MAP1, randomize known stimulus coordinates at sensor centres, between
  sensors, and near boundaries. The operator viewing the software is blinded
  to the true coordinate. Save the stimulus coordinate, detected coordinate,
  nearest-neighbour pitch, interpolation method, coverage mask, and confidence
  flag.\par
  \[
  e_{loc}=\sqrt{(x_d-x_t)^2+(y_d-y_t)^2}
  \]
  Report median, 95th percentile, and boundary error; do not report only a
  visually convincing heatmap.\par}
\end{minipage}

{\fontsize{7.45}{8.6}\selectfont\color{AUBGray}
\textbf{Figure Q5-5. Blinded localization test.} Blue dots are measured sensor
sites, red is the applied coordinate, and yellow is the algorithm output. The
Q4 gate is median error no greater than half the nearest-neighbour pitch;
sensor markers and unsupported regions remain visible.\par}

\vspace{0.35em}
\evidencebox{Expected outcome and impact:}{A successful run preserves channel
identity, unit, timestamp, calibration version, and CAD coordinate; localizes
held-out inputs within the Q4 gate; and leaves pressure, signed shear,
temperature, and RH viewable separately. Failure changes sampling,
synchronization, shielding, channel compensation, sensor density, coordinate
registration, interpolation, or the displayed claim. The map remains an
\emph{interface-condition map} until participant evidence validates a
discomfort or clinical-risk interpretation.}

% -----------------------------------------------------------------------------
\clearpage
\hypertarget{q5-lit-integration}{}
\sectiontitle{Q5 Literature Evidence Plate: Multimodal Mapping}

{\fontsize{8.75}{10.25}\selectfont
The literature already provides reasonable solutions for synchronized
multimodal acquisition, embedded liner microclimate sensing, and 3D surface
visualization. X1 and MAP1 therefore test whether those functional ideas remain
valid for the present printed-slab geometry and CAD manifest; see
\reportlink{q5-integration-mapping}{Q5 Section 5.5} and
\reportlink{q3-measurement-models}{Q3 Section 3.5}.\par}

\vspace{0.55em}
\noindent
\qfivelitpanel{0.255\textheight}
  {paper_sources/q5_figures/crops/kwak2020_fig1.png}
  {\textbf{(a) Wireless multimodal prosthetic interface.} Original Fig.~1 in
  Kwak \textit{et al.}: pressure and temperature sensors, NFC/BLE modules,
  wearable placement, and the functional data path \cite{q3kwak2020}.
  \href{https://doi.org/10.1126/scitranslmed.abc4327}{Open DOI}. The figure is
  a system precedent, not evidence for the present silver/PDMS channels.}
\hfill
\qfivelitpanel{0.255\textheight}
  {paper_sources/q5_figures/crops/turner2021_fig3.png}
  {\textbf{(b) Interpolated socket-surface display.} Original Fig.~3 in
  Turner \textit{et al.}: radial-basis interpolation applied to a 3D
  transtibial socket model \cite{q3turner2021}. It supports a live
  coordinate-registered visualization while reinforcing the need to report
  true-versus-detected location error.
  \href{https://doi.org/10.3390/prosthesis3040035}{Open DOI}.}

\vspace{0.75em}
\begin{center}
\includegraphics[width=0.88\linewidth,height=0.25\textheight,keepaspectratio]
  {paper_sources/figures_selected/paterno_fig4_temperature_humidity_results.png}
\end{center}
{\fontsize{7.0}{8.05}\selectfont
\textbf{(c) Embedded-liner microclimate data.} Original Fig.~4 in Patern\`o
\textit{et al.}: thermography, sensor locations, and temperature/RH traces
from a personalised liner with embedded sensors \cite{paterno2020}. It is a
direct precedent for synchronized physical-unit channels and location-aware
interpretation. \href{https://doi.org/10.1186/s12938-020-00814-y}{Open DOI}.\par}

\vfill
{\fontsize{6.7}{7.7}\selectfont\color{AUBGray}
\textbf{Figure Q5-L4. Existing functional solutions leveraged by X1/MAP1.}
The novel project question is not whether maps or multimodal loggers can exist;
it is whether calibrated signals from the present sensor slabs retain channel,
time, and coordinate identity well enough to support an accurate map.\par}

% -----------------------------------------------------------------------------
\clearpage
\hypertarget{q5-durability}{}
\sectiontitle{Q5. Environmental, Cyclic, and Recovery Testing}

\qtwosubsection{5.6 D1--D2 --- measure drift instead of assuming durability}

{\fontsize{8.9}{10.45}\selectfont
Directly printed silver conductors on PDMS have shown strong dependence on
strain, recovery time, and cycle count; Albrecht \textit{et al.} reported
conductive traces after 10,000 cycles but with a large resistance increase
under their specific geometry and 20\% strain protocol
\cite{albrecht2018stretch}. Kwak \textit{et al.} reported 1,000-cycle soft
prosthetic-sensor testing, while Patern\`o \textit{et al.} demonstrated
temperature/RH operation in a liner microclimate \cite{q3kwak2020,paterno2020}.
These are durability precedents, not substitutes for the present slab's
post-exposure recalibration. Q5 follows the release gate in
\reportlink{q4-design-gate}{Q4 Section 4.7}.\par}

\begin{center}
\begin{tikzpicture}[>=stealth,node distance=2.7mm]
  \tikzstyle{dur}=[draw=AUBLine,fill=AUBPale,rounded corners=1pt,
    text width=30mm,minimum height=10mm,align=center,
    font=\fontsize{6.6}{7.5}\selectfont]
  \node[dur] (a) {baseline F0 + modality calibration};
  \node[dur,right=of a] (b) {35--37~$^{\circ}$C + high RH conditioning};
  \node[dur,right=of b] (c) {artificial sweat + protected exposure};
  \node[dur,right=of c] (d) {1,000 then 10,000 load/bend cycles};
  \node[dur,right=of d] (e) {24 h recovery + full recalibration};
  \draw[->,thick,color=AUBRed] (a)--(b);
  \draw[->,thick,color=AUBRed] (b)--(c);
  \draw[->,thick,color=AUBRed] (c)--(d);
  \draw[->,thick,color=AUBRed] (d)--(e);
\end{tikzpicture}
\end{center}
{\fontsize{7.5}{8.7}\selectfont\color{AUBGray}
\textbf{Figure Q5-6. D1--D2 exposure sequence.} Intermediate inspections at
defined cycle counts distinguish progressive damage from a single final
failure. Conditioning, load, and recovery records share the specimen ID.\par}

\vspace{0.4em}
\begingroup
\fontsize{7.65}{8.95}\selectfont
\setlength{\tabcolsep}{2.55pt}\renewcommand{\arraystretch}{1.04}
\begin{tabularx}{\linewidth}{>{\raggedright\arraybackslash\bfseries\color{AUBRed}}p{1.05cm}Y Y Y}
\rowcolor{AUBRed}
\color{white}\textbf{Test} & \color{white}\textbf{Exposure and tools} &
\color{white}\textbf{Measurements} & \color{white}\textbf{Expected outcome and impact} \\
D1-1 & Chamber or sealed enclosure at 35--37~$^{\circ}$C and high RH; reference
temperature/RH; powered slab where safe. & Baseline, resistance/capacitance,
noise, leakage, drift, visual condition, and modality calibration before,
during, and after conditioning. & No damage and no Q4-exceeding calibration
shift after recovery. Failure changes materials, encapsulation, compensation,
operating range, or recalibration interval. \\
\rowcolor{AUBPale}
D1-2 & Controlled artificial sweat, selected barrier, absorbent containment,
camera, electrical isolation/leakage checks. & Liquid path, saturation,
recovery, open/short status, adhesion, corrosion/discolouration, and H1
recalibration. & No liquid contact with exposed silver or joints; vapour
response recovers. Failure changes barrier, seal, vent position, replaceable
layer, or cleaning instructions. \\
D2-1 & Cyclic normal compression and signed shear fixture; force/displacement
reference; staged 1,000-cycle screen then 10,000-cycle characterization. & Raw
signals at defined cycle intervals; $\Delta R/R_0$, $\Delta C/C_0$, sensitivity,
hysteresis, repeatability, baseline recovery, and inspection images. & No
open/short, delamination, or unsafe exposure; performance remains inside Q4
gates. Failure changes trace geometry, neutral-axis placement, PDMS stack,
bonding, or claimed service interval. \\
\rowcolor{AUBPale}
D2-2 & Mandrels or curved phantom representing intended liner curvature;
controlled bend radius and cycle count. & Continuity during bending,
strain-only cross-sensitivity, crack location, and post-bend modality
calibration. & Bending must not be interpreted as temperature, humidity, or
load without compensation. Failure changes routing, strain relief, sensor
orientation, or enclosure geometry. \\
\end{tabularx}
\endgroup

\vspace{0.35em}
\begin{align*}
D_y(\%)&=\frac{y_{post}-y_{pre}}{FS}\times100, &
\Delta R/R_0&=\frac{R_N-R_0}{R_0}, &
CV(\%)&=100\,s/\bar y.
\end{align*}
{\fontsize{7.5}{8.7}\selectfont\color{AUBGray}
Report raw and normalized changes because a small percentage change near zero
can be misleading. Recovery at 5, 30, 300~s and 24~h separates short-term PDMS
relaxation from persistent damage.\par}

% -----------------------------------------------------------------------------
\clearpage
\hypertarget{q5-phantom}{}
\sectiontitle{Q5. Simulated-Use Phantom Testing}

\qtwosubsection{5.7 S1 --- controlled curved-interface scenarios before human use}

{\fontsize{8.9}{10.45}\selectfont
The three build precedents, photographs, and material lists are retained in
\reportlink{q3-phantom-precedents}{Q3 Section 3.4.1} and are not duplicated
here. Patern\`o \textit{et al.} built an anatomically informed variable-volume
transfemoral simulator using silicone tissue layers, a sealed hydraulic
chamber, and motorized syringe pumps \cite{paterno2025socket}. McGrath
\textit{et al.} built a silicone residuum with molded pores that released
water inside a liner/socket under cyclic loading \cite{mcgrath2021sweat}.
Phillips \textit{et al.} reported a lower-cost dual-hardness urethane mock limb
with water-filled TPU bladders and motorized syringe control
\cite{phillips2025}. Q5 uses these as selectable modules because no single
phantom provides traceable normal load, signed shear, vapour RH, liquid sweat,
temperature, and anatomical motion simultaneously.\par}

\begingroup
\fontsize{7.55}{8.85}\selectfont
\setlength{\tabcolsep}{2.5pt}\renewcommand{\arraystretch}{1.04}
\begin{tabularx}{\linewidth}{>{\raggedright\arraybackslash\bfseries\color{AUBRed}}p{2.0cm}Y Y Y}
\rowcolor{AUBRed}
\color{white}\textbf{Candidate module} & \color{white}\textbf{Materials and tools} &
\color{white}\textbf{Q5 role and ground truth} &
\color{white}\textbf{Limitation and decision impact} \\
A --- variable-volume anatomical simulator & Molded silicone anatomical/tissue
layers, sealed hydraulic chamber, syringe pumps, frame, volume/displacement
reference, force/load references. & Apply repeatable volume changes and
interface loading while checking pressure/shear localization, coverage, and
map stability. & Higher fabrication complexity; it does not by itself provide
controlled vapour RH. Select when anatomical loading and volume change are the
dominant questions. \\
\rowcolor{AUBPale}
B --- sweating residuum/socket simulator & 3D-printed bone/core, silicone,
molded pores, water reservoir/pump, liner/check socket, cyclic actuator,
collected-volume measurement. & Challenge liquid-sweat protection and
moisture-induced changes during cyclic loading; record delivered and recovered
water volume. & Liquid delivery is not RH calibration. Use only after H1 dry
vapour calibration; select when sweat management and wet slip are priorities. \\
C --- low-cost bladder mock limb & Dual-hardness urethane, water-filled TPU
bladders, mold/printed parts, Arduino-class control, motorized syringe, force/
displacement reference. & Reproduce controlled shape/pressure change at lower
reported construction cost and compare known actuation with mapped response. &
Biofidelity and shear require separate validation. Select when reproducible
low-cost pressure/shape testing is the priority. \\
\end{tabularx}
\endgroup

\vspace{0.45em}
\begin{center}
\begin{tikzpicture}[>=stealth,node distance=2.6mm]
  \tikzstyle{sim}=[draw=AUBLine,fill=AUBPale,rounded corners=1pt,
    text width=31mm,minimum height=9mm,align=center,
    font=\fontsize{6.6}{7.5}\selectfont]
  \node[sim] (a) {dry seated baseline\\known coordinate zero};
  \node[sim,right=of a] (b) {known normal load / volume change};
  \node[sim,right=of b] (c) {signed motion or cyclic loading};
  \node[sim,right=of c] (d) {controlled liquid delivery if Module B};
  \node[sim,right=of d] (e) {unload, dry recovery, recalibrate};
  \draw[->,thick,color=AUBRed] (a)--(b);
  \draw[->,thick,color=AUBRed] (b)--(c);
  \draw[->,thick,color=AUBRed] (c)--(d);
  \draw[->,thick,color=AUBRed] (d)--(e);
\end{tikzpicture}
\end{center}
{\fontsize{7.45}{8.6}\selectfont\color{AUBGray}
\textbf{Figure Q5-7. S1 simulated-use sequence.} Ground truth comes from the
phantom actuator, reference load/displacement, delivered liquid volume, and
known CAD locations. The sensor map is the device under test, not the source of
the event label.\par}

\vspace{0.4em}
\evidencebox{Expected outcome and impact:}{The curved assembly retains
electrical continuity and coordinate identity, detects the correct imposed
event and location within the Q4 gates, and returns toward baseline after
unloading/drying. Failure changes slab placement, liner integration, routing,
encapsulation, protection, calibration, interpolation, or the chosen phantom
module. This is simulated-use evidence only; it does not authorize a human-use
or clinical discomfort claim.}

% -----------------------------------------------------------------------------
\clearpage
\hypertarget{q5-lit-phantoms}{}
\sectiontitle{Q5 Literature Evidence Plate: Residual-Limb Phantoms}

{\fontsize{8.75}{10.25}\selectfont
The three proposed modules in \reportlink{q5-phantom}{Q5 Section 5.7} are all
implemented research systems. Their published photographs are shown together
so the project can choose a build route from demonstrated hardware rather than
from a generic sketch. Their separate strengths also explain why no one module
is treated as a universal ground truth.\par}

\vspace{0.55em}
\noindent
\qfivelitpanel{0.245\textheight}
  {assets/q3/phantoms/paterno2025_fig5a.png}
  {\textbf{(a) Anatomical variable-volume simulator.} Original Fig.~5a in
  Patern\`o \textit{et al.}: anatomical skeletal/tissue structure, silicone
  layers, sealed hydraulic chamber, and two motorized syringe pumps for
  controlled residual-limb volume change \cite{paterno2025socket}.
  \href{https://doi.org/10.1002/aisy.202400995}{Open DOI}.}
\hfill
\qfivelitpanel{0.245\textheight}
  {assets/q3/phantoms/mcgrath2021_fig1.jpg}
  {\textbf{(b) Sweating residuum/socket simulator.} Original Fig.~1 in
  McGrath \textit{et al.}: a 3D-printed bone/core encased in porous silicone,
  supplied with water and tested inside a liner/check socket under cyclic
  loading \cite{mcgrath2021sweat}.
  \href{https://doi.org/10.33137/cpoj.v4i1.35213}{Open DOI}.}

\vspace{0.8em}
\begin{center}
\includegraphics[width=0.83\linewidth,height=0.25\textheight,keepaspectratio]
  {assets/q3/phantoms/phillips2025_fig1.png}
\end{center}
{\fontsize{7.0}{8.05}\selectfont
\textbf{(c) Lower-cost biofidelic mock residual limb.} Original Fig.~1 in
Phillips \textit{et al.}: dual-hardness urethane tissues, water-filled TPU
bladders, mold components, and a motorized syringe/Arduino-class control route
reported at below CAD~400 \cite{phillips2025}.
\href{https://doi.org/10.33137/cpoj.v8i2.45759}{Open DOI}.\par}

\vfill
{\fontsize{6.7}{7.7}\selectfont\color{AUBGray}
\textbf{Figure Q5-L5. Implemented phantom candidates.} Patern\`o supplies
anatomical volume change, McGrath supplies controlled liquid sweat during
cyclic loading, and Phillips supplies a reproducible lower-cost pressure/shape
platform. These images are duplicated here intentionally because Q5 uses them
to choose test hardware; the detailed materials comparison remains in
\reportlink{q3-phantom-precedents}{Q3 Section 3.4.1}.\par}

% -----------------------------------------------------------------------------
\clearpage
\hypertarget{q5-reporting}{}
\sectiontitle{Q5. Analysis Outputs, Decision Rules, and Design Feedback}

\qtwosubsection{5.8 SYS1 and R1 --- convert each test into auditable evidence}

{\fontsize{9.0}{10.6}\selectfont
Every test produces a raw file, reference file, specimen/build record, analysis
script or calculation sheet, figure, numerical metric table, and decision.
The output formats below implement the metrics defined in
\reportlink{q3-metrics}{Q3 Section 3.7} and the planned evidence panels in
\reportlink{q3-validation-outputs}{Q3 Section 3.8}. Literature figures may
explain precedent, but only project records occupy the measured-result panels.\par}

\begingroup
\fontsize{7.55}{8.85}\selectfont
\setlength{\tabcolsep}{2.55pt}\renewcommand{\arraystretch}{1.04}
\begin{tabularx}{\linewidth}{>{\raggedright\arraybackslash\bfseries\color{AUBRed}}p{1.05cm}Y Y Y}
\rowcolor{AUBRed}
\color{white}\textbf{Output} & \color{white}\textbf{Required plot/table} &
\color{white}\textbf{Decision rule} & \color{white}\textbf{Design feedback} \\
F0 & Annotated print photograph, CAD-versus-measured dimensions, continuity/
isolation table, build traveller. & Advance only intact, identified, traceable
specimens; otherwise rework/reject with cause. & Print parameters, mold,
surface preparation, alignment, contacts, or encapsulation. \\
\rowcolor{AUBPale}
M1--M3 & Reference input and raw channels versus time; calibration with
residuals; loading/unloading hysteresis; cross-axis and recovery; localization
error map. & Apply MECH-01--MECH-05 individually. Averages cannot hide a
failed direction, boundary, or recovery gate. & Electrode/spacer geometry,
preload, footprint, placement, compensation, or coverage claim. \\
T1/H1/H2 & Reference and sensor versus time; fitted calibration and held-out
error; $t_{90}$; hysteresis; temperature compensation; pre/post-liquid shift.
& Apply TEMP-01/02 and HUM-01--03 separately. Saturation or quality failure is
flagged rather than converted into a condition score. & Meander/IDE geometry,
barrier, vent, seal, compensation, range, or recalibration. \\
\rowcolor{AUBPale}
X1/MAP1 & Modality-by-modality matrix; timestamp audit; true/detected coordinate
scatter; error vectors; coverage/confidence mask; separate physical-unit maps.
& Apply DAQ-01/02, MAP-01--03, and ALG-01. Combined interpretation is blocked
when calibration, time, coordinate, or quality metadata are missing. & DAQ,
shielding, synchronization, manifest, sensor density, interpolation, thresholds,
or displayed wording. \\
D1/D2/S1 & Drift/recovery versus exposure or cycle count; inspection images;
pre/post calibration; event and location results against phantom ground truth.
& Apply DUR, ENV, SAFE, and SYS gates per modality and specimen. Report
pass/partial/fail/not run with the evidence path. & Materials, routing, barrier,
stack, service interval, phantom module, or decision to remain at flat-slab
testing. \\
R1 & Batch yield, quantities, supplier/unit cost, reusable equipment, labour,
failed-print causes, and AUB availability status. & No guessed availability or
cost; attach inventory confirmation, receipts/quotes, and build time. & Process
simplification, supplier choice, sensor count, reusable fixtures, or scope. \\
\end{tabularx}
\endgroup

\vspace{0.45em}
\begin{center}
\begin{tikzpicture}[x=1cm,y=1cm]
  \draw[fill=AUBPale,draw=AUBLine] (0,0) rectangle (4.0,1.15);
  \draw[fill=AUBPale,draw=AUBLine] (4.3,0) rectangle (8.3,1.15);
  \draw[fill=AUBPale,draw=AUBLine] (8.6,0) rectangle (12.6,1.15);
  \draw[fill=AUBPale,draw=AUBLine] (12.9,0) rectangle (16.9,1.15);
  \node[font=\fontsize{7}{8}\selectfont,align=center] at (2,0.72)
    {\textbf{PASS}\\all applicable gates met};
  \node[font=\fontsize{7}{8}\selectfont,align=center] at (6.3,0.72)
    {\textbf{PARTIAL}\\usable block; stated limitation};
  \node[font=\fontsize{7}{8}\selectfont,align=center] at (10.6,0.72)
    {\textbf{FAIL}\\requirement not met; redesign};
  \node[font=\fontsize{7}{8}\selectfont,align=center] at (14.9,0.72)
    {\textbf{NOT RUN}\\no conclusion; reason recorded};
  \draw[AUBRed,thick] (0,-0.18)--(16.9,-0.18);
\end{tikzpicture}
\end{center}
{\fontsize{7.5}{8.7}\selectfont\color{AUBGray}
\textbf{Figure Q5-8. Required verdict vocabulary.} A partial result preserves a
useful functional block while identifying its limitation; ``not run'' is not a
failure and cannot be presented as evidence.\par}

\vspace{0.45em}
\evidencebox{Q5 completion criterion:}{Q5 is complete when every planned test
has a specimen ID, apparatus/reference description, controlled inputs, raw and
processed outputs, metric definitions, acceptance rule, actual status, and
design action. Before data collection, the section is correctly reported as a
literature-backed validation plan. After testing, planned panels are replaced
with measured project plots without changing the preregistered decision rules.}

\sourceanchor{Bench methods and performance reporting are grounded in printed
force, temperature, humidity, PDMS-conductor, prosthetic-interface, mapping,
and phantom studies
\cite{albrecht2019,liew2020,jager2022,aitmammar2020,
albrecht2018stretch,q3kwak2020,paterno2020,q3turner2021,
paterno2025socket,mcgrath2021sweat,phillips2025}.}
